\subsection{关于连续线性表示的回顾}

设 G 为拓扑群。我们回忆，G 在拓扑 K-向量空间 E 中的连续线性表示是 G 在 E 中的线性表示 \(\varrho\) ，使得由 \((g, x) \mapsto \varrho(g)x\) 定义的从 G × E 到 E 的映射连续（INT, VIII, 第 128 页, § 2, n° 1, def. 1）。这个映射定义了 G 在 E 上的一个作用，且 \(\varrho\) 的像包含于 \(\mathcal{L}(\mathrm{E})\) 。我们称连续表示 \(\varrho\) 是有界的，如果它的像在赋有有界收敛拓扑的空间 \(\mathcal{L}(\mathrm{E})\) 中有界。

注. —— 若 K = R 且 \(\varrho\) 是 G 在 R-向量空间 E 中的连续线性表示，则映射 \(\varrho_{(\mathbf{C})}: g \mapsto 1 \otimes \varrho(g)\) 是 G 在 \(\mathrm{E}_{(\mathbf{C})}\) 中的连续线性表示；若 \(\varrho\) 有界，则 \(\varrho_{(\mathbf{C})}\) 有界。

对每个拓扑 K-向量空间 E，把每个 \(g \in G\) 映为 \(1_{E}\) 的同态是 G 在 E 中的连续线性表示，称为 G 在 E 中的平凡表示。

设 H 为 G 的子群。\(\varrho\) 到 H 的限制是 H 在 E 中的连续线性表示，记作 \(\mathrm{Res}_{\mathrm{H}}^{\mathrm{G}}(\varrho)\) 。

设 \(\varrho_{1}\) 与 \(\varrho_{2}\) 为 G 在拓扑 K-向量空间 E \(_{1}\) 与 E \(_{2}\) 中的连续线性表示。G-态射 \(u\) （从 \(\varrho_{1}\) 到 \(\varrho_{2}\) ）是连续的线性表示态射，即连续线性映射 \(u\colon\) E \(_{1} \to\) E \(_{2}\) ，满足 \(u \circ \varrho_{1}(g) = \varrho_{2}(g) \circ u\) （对所有 \(g \in\) G）。我们用 Hom \(_{G}\) ( \(\varrho_{1}, \varrho_{2}\) ) 表示从 \(\varrho_{1}\) 到 \(\varrho_{2}\) 的 G-态射组成的向量空间。它是赋有逐点收敛拓扑的空间 \(\mathcal{L}\) (E \(_{1}\) ; E \(_{2}\) ) 的闭子空间。

设 \(\varrho\) 为群 G 在拓扑 K-向量空间 E 中的连续线性表示。E 的恒等映射 \(1_{\mathrm{E}}\) 是 \(\varrho\) 到 \(\varrho\) 的 G-态射，记作 \(1_{\varrho}\) 。若 \(\varrho_1\) 、 \(\varrho_2\) 与 \(\varrho_3\) 分别是 G 在拓扑 K-向量空间 \(\mathrm{E}_1\) 、 \(\mathrm{E}_2\) 、 \(\mathrm{E}_3\) 中的连续线性表示，且 \(u: \mathrm{E}_1 \to \mathrm{E}_2\) 与 \(v: \mathrm{E}_2 \to \mathrm{E}_3\) 是 G-态射，则 \(v \circ u\) 是 G-态射。

G-态射 \(u\) （从 \(\varrho_{1}\) 到 \(\varrho_{2}\) ）是 G-同构，如果存在 G-态射 \(v\) （从 \(\varrho_{2}\) 到 \(\varrho_{1}\) ），使 \(v\circ u=1_{\varrho_{1}}\) 且 \(u\circ v=1_{\varrho_{2}}\) 。为此必要且充分的是：\(u\) 是 G-态射，且是空间 \(\varrho_{1}\) 到空间 \(\varrho_{2}\) 的拓扑向量空间同构；其逆 \(v=u^{-1}\) 那时确实是 G-态射。若存在从 \(\varrho_{1}\) 到 \(\varrho_{2}\) 的 G-同构，我们就说这些表示是同构的。

定义 1.—— 设 \(\varrho\) 为拓扑群 G 在拓扑向量空间 E 中的连续线性表示。我们称 \(\pi\) 是 \(\varrho\) 的一个子表示，如果 F 是 E 的闭子空间，且对每个 \(g \in G\) ，空间 F 在 \(\varrho(g)\) 下稳定，而 \(\pi(g)\) 是 \(\varrho(g)\) 通过过渡到子空间诱导的 F 的自同态。

设 \(\varrho\) 为拓扑群 G 在拓扑向量空间 E 中的连续线性表示。\(\{0\}\) 在 E 中的闭包是 \(\varrho\) 的子表示。更一般地，在诸自同态 \(\varrho(g)\) 下稳定的子空间的闭包是 \(\varrho\) 的子表示。

子表示 \(\pi\) （属于 \(\varrho\) ）由一个在所有自同态 \(\varrho(g)\) 下稳定的闭子空间 F 唯一确定。后一条件常表述为“F 是 E 的 G-不变子空间”。这时我们也说 F 定义了 \(\varrho\) 的一个子表示，或者有时滥用语言，说 F 是 \(\varrho\) 或 E 的子表示。

设 F 为 E 的一个子空间，它定义了子表示 \(\pi\) （\(\varrho\) 的）。对每个 \(g \in G\) ，线性映射 \(\varrho(g)\) 通过过渡到商定义了 E/F 的一个自同态 \(\widetilde{\varrho}(g)\) 。映射 \(g \mapsto \widetilde{\varrho}(g)\) 是 G 在 E/F 中的连续线性表示，且 E 到 E/F 的典范投影是 G-态射。我们称 \(\widetilde{\varrho}\) 为 G 在 E/F 上的商表示；它也记作 \(\varrho/\pi\) 。

E 中在 G 的作用下不变的元素组成的子空间是 \(\varrho\) 的平凡子表示。我们把它记作 \(\varrho^{G}\) 或 \(E^{G}\) 。

设 A 为 E 的子集。由元素 \(\varrho(g)x\) （其中 \(g \in G\) ，\(x \in A\) ）生成的闭子空间 F 是 \(\varrho\) 的子表示，称为由 A 生成的 \(\varrho\) 的子表示；若 F = E，就说 A 生成 \(\varrho\) 。若子表示 F 是有限维的，就说子集 A 是 G-有限的。

假设 A 缩减为单个元素 \(x \in E\) 。若 A 生成 \(\varrho\) ，就称 x 是 \(\varrho\) 的循环向量，并称 \(\varrho\) 是循环表示；若 A 是 G-有限的，就称 x 是 G-有限向量。

\(\varrho\) 的 G-有限向量组成的集是 E 的在 \(\varrho\) 下稳定的向量子空间；它不一定是闭的。

设 \((\varrho_{i})_{i\in\mathrm{I}}\) 为 G 在局部凸拓扑 K-向量空间 \((\mathrm{E}_{i})_{i\in\mathrm{I}}\) 中的一族连续线性表示。设 E 为诸空间 \(E_{i}\) 的直和空间。映射 \(g\mapsto(\varrho_{i}(g))\) 是 G 在 E 中的线性表示，称为诸表示 \(\varrho_{i}\) 的和或直和。若 I 有限，则它是连续的。

若表示 \(\varrho_{i}\) 都等于某个表示 \(\varrho\) （对所有 \(i \in I\) ），我们就说诸表示 \(\varrho_{i}\) 的直和是表示 \(\varrho\) 的 Card(I) 个拷贝之和，它也记作 \(\varrho^{\mathrm{Card(I)}}\) ，或 Card(I) \(\varrho\) 。

设 \((\varrho_{i})_{i\in\mathrm{I}}\) 与 \((\pi_{j})_{j\in\mathrm{J}}\) 分别为 G 在拓扑 K-向量空间 \((\mathrm{E}_{i})_{i\in\mathrm{I}}\) 与 \((\mathrm{F}_{j})_{j\in\mathrm{J}}\) 中的有限族连续线性表示。K-向量空间的典范同构

\[
\operatorname{Hom} _ {\mathrm{K}} \left(\bigoplus_ {i \in \mathrm{I}} \mathrm{E} _ {i}, \bigoplus_ {j \in \mathrm{J}} \mathrm{F} _ {j}\right)\rightarrow \bigoplus_ {(i, j) \in \mathrm{I} \times \mathrm{J}} \operatorname{Hom} _ {\mathrm{K}} (\mathrm{E} _ {i}, \mathrm{F} _ {j})
\]

（A, II, 第 13 页, cor. 1）通过过渡到子空间诱导出同构

\[
\mathrm{Hom} _ {\mathrm{G}} \Bigl (\bigoplus_ {i \in \mathrm{I}} \varrho_ {i}, \bigoplus_ {j \in \mathrm{J}} \pi_ {j} \Bigr) \to \bigoplus_ {(i, j) \in \mathrm{I} \times \mathrm{J}} \mathrm{Hom} _ {\mathrm{G}} (\varrho_ {i}, \pi_ {j})\tag{1}
\]

同样称它是典范的。

设 \(\varrho_{1}\) 与 \(\varrho_{2}\) 为 G 在 K-向量空间 E \(_{1}\) 与 E \(_{2}\) 中的线性表示。对 \(u \in \mathcal{L}(\mathrm{E}_{1};\mathrm{E}_{2})\) 与 \(g \in \mathrm{G}\) ，我们令 \(\varrho(g)u = \varrho_{2}(g) \circ u \circ \varrho_{1}(g^{-1})\) 。映射 \(g \mapsto \varrho(g)\) 是 G 在 \(\mathcal{L}(\mathrm{E}_{1};\mathrm{E}_{2})\) 中的线性表示。这个表示的不变元素组成的空间与 \(\mathrm{Hom}_{\mathrm{G}}(\varrho_{1},\varrho_{2})\) 重合。

设 \(\varrho\) 为 G 在局部凸空间 E 中的连续线性表示。回忆（INT, VIII, 第 131 页, § 2, no. 2）逆步表示 \(\breve{\varrho}\) （\(\varrho\) 的）是 G 在 E 的对偶 E' 中的线性表示，由 \(\breve{\varrho}(g) = {}^{t} \varrho(g^{-1})\) 定义。

